\documentclass[12pt]{article}

\usepackage{sbc-template}
\usepackage{graphicx,url}
\usepackage[utf8]{inputenc}
\usepackage[brazil]{babel}

\sloppy

\title{Simulação Concorrente de um Centro de\\ Distribuição Automatizado (IDP)}

\author{Samuel Abrão\inst{1}, Andreia Tiveron\inst{1}, João Genaro\inst{1}}

\address{Instituto Brasileiro de Ensino, Desenvolvimento e Pesquisa (IDP)\\
  Asa Norte -- Brasília -- DF -- Brasil
  \email{samuelabrao2006@gmail.com, andreiabtiveron@gmail.com,}\\
  \texttt{\footnotesize genastrb@gmail.com}}

\begin{document}

\maketitle

\begin{abstract}
This report describes a concurrent simulation, written in C with POSIX threads,
of an automated distribution center. Collector and deliverer robots share a
grid map, pickup stations and a bounded conveyor belt, moving one cell at a time
under the physical constraint that two entities never occupy the same cell. We
detail the thread organization, the shared resources, the identified critical
sections, the synchronization mechanisms (fine-grained mutexes and a condition
variable) and the main design decisions, followed by the build and run guide for
Linux (Ubuntu 24.04).
\end{abstract}

\begin{resumo}
Este documento descreve uma simulação concorrente, escrita em C com \textit{threads}
POSIX, de um centro de distribuição automatizado. Robôs coletores e entregadores
compartilham um mapa em grade, estações de coleta e uma esteira transportadora
limitada, movendo-se célula a célula sob a restrição física de que duas entidades
nunca ocupam a mesma célula. Detalham-se a organização das \textit{threads}, os
recursos compartilhados, as seções críticas identificadas, os mecanismos de
sincronização (\textit{mutexes} de granularidade fina e uma variável de condição)
e as principais decisões de projeto, seguidos do guia de compilação e execução em
Linux (Ubuntu 24.04).
\end{resumo}

\section{Introdução}

A simulação reproduz o centro de distribuição da IDP, dividido em três regiões:
a \textbf{área de coleta}, com estações geradoras de pacotes (pontos \texttt{P});
a \textbf{esteira transportadora}, ligando a entrada (\texttt{in}) à saída
(\texttt{out}); e a \textbf{área de expedição}, com pontos de despacho
(\texttt{D}). O fluxo de um pacote é: o \emph{gerador} o cria em uma estação
\texttt{P}; um \emph{robô coletor} o transporta até \texttt{in}; a \emph{esteira}
o desloca uma posição por unidade de tempo até \texttt{out}; e um \emph{robô
entregador} o retira e o leva a um ponto \texttt{D}. A restrição física central é
que \emph{duas entidades nunca ocupam a mesma célula do mapa}. A simulação termina
quando a quantidade de pacotes do cenário é entregue, ou quando o usuário pressiona
\texttt{q}.

O sistema foi escrito em C com \texttt{pthreads}~\cite{tanenbaum} e interface em
\texttt{ncurses}~\cite{ncurses}. Seguindo a recomendação do enunciado, primeiro
foi construída e validada uma versão sequencial e só então introduziu-se a
concorrência. O código é organizado em módulos independentes que correspondem aos
elementos exigidos: \texttt{gerador} (Gerador de Pacotes), \texttt{coletor}
(Robôs Coletores), \texttt{esteira} (Esteira Transportadora), \texttt{entregador}
(Robôs Entregadores) e \texttt{interface} (Interface); os módulos \texttt{mapa} e
\texttt{estacao} implementam os recursos compartilhados de base. O mapa é
representado por duas matrizes: uma de tipos de célula (livre, parede, estação
\texttt{P}, ponto \texttt{D}) e uma de ocupação; cada pacote carrega um estado
(\emph{aguardando}, \emph{transportado}, \emph{na esteira}, \emph{entregue}).
São oferecidos três cenários estáticos, selecionáveis por argumento de linha de
comando ou por menu (Tabela~\ref{tab:cenarios}).

\begin{table}[ht]
\centering
\caption{Parâmetros dos três cenários estáticos.}
\label{tab:cenarios}
\small
\begin{tabular}{lccc}
\hline
\textbf{Parâmetro} & \textbf{0 -- Básico} & \textbf{1 -- Intermediário} & \textbf{2 -- Estresse} \\
\hline
Mapa (largura $\times$ altura) & $20\times10$ & $30\times15$ & $40\times20$ \\
Robôs coletores        & 2 & 4 & 6 \\
Robôs entregadores     & 2 & 3 & 5 \\
Estações \texttt{P}    & 2 & 3 & 4 \\
Pontos \texttt{D}      & 2 & 3 & 4 \\
Tamanho da esteira     & 8 & 12 & 20 \\
Total de pacotes       & 20 & 50 & 100 \\
\hline
\end{tabular}
\end{table}

\section{Organização das Threads}

Cada entidade da simulação é uma \textit{thread} POSIX independente, criada com
\texttt{pthread\_create} e finalizada com \texttt{pthread\_join} (nenhuma
\textit{thread} é \textit{detached}). A \textit{thread} principal não modela uma
entidade: ela desenha a interface (é a única a chamar \texttt{ncurses}),
acompanha o critério de término e faz o \texttt{join} das demais. Cada
\textit{thread} de entidade executa um laço que dorme por um passo de tempo
(\texttt{TICK\_US}, 100\,ms; coletores e entregadores usam metade do intervalo) e
verifica, a cada iteração, uma \textit{flag} global de encerramento.

\begin{itemize}
  \item \textbf{Gerador} (1 \textit{thread}): a cada passo cria um pacote em uma
        estação \texttt{P}, em \textit{round-robin} com \textit{liveness} (busca
        outra estação com espaço se a atual estiver cheia);
  \item \textbf{Esteira} (1 \textit{thread}): a cada passo avança os pacotes uma
        posição rumo a \texttt{out}, se a posição seguinte estiver livre;
  \item \textbf{Coletores} ($N$ \textit{threads}): autômato de três estados
        --- \texttt{OCIOSO} (procura estação com pacote), \texttt{INDO\_ESTACAO}
        (desloca-se e coleta) e \texttt{INDO\_ESTEIRA} (leva o pacote até
        \texttt{in} e o insere; se \texttt{in} estiver ocupada, aguarda);
  \item \textbf{Entregadores} ($M$ \textit{threads}): autômato de três estados
        --- \texttt{OCIOSO} (espera pacote em \texttt{out}), \texttt{INDO\_ESTEIRA}
        (vai a \texttt{out} e retira) e \texttt{INDO\_DESPACHO} (leva o pacote a
        um ponto \texttt{D} e o contabiliza).
\end{itemize}

O total é $2 + N + M$ \textit{threads} de entidades mais a principal (por exemplo,
seis \textit{threads} de entidades no cenário básico). O encerramento é
coordenado: uma \textit{flag} booleana global, protegida por \texttt{mutex}, é
consultada por cada \textit{thread} a cada iteração e sinalizada pela principal
quando todos os pacotes são entregues ou o usuário tecla \texttt{q}; em seguida,
todas as \textit{threads} são unidas com \texttt{pthread\_join}.

\section{Recursos Compartilhados}

Cada recurso possui trava própria, de granularidade fina, em vez de um único
\textit{lock} global, permitindo progresso paralelo. A Tabela~\ref{tab:recursos}
resume o mapeamento recurso $\rightarrow$ mecanismo.

\begin{table}[ht]
\centering
\caption{Recursos compartilhados e mecanismos de sincronização.}
\label{tab:recursos}
\footnotesize
\setlength{\tabcolsep}{4pt}
\begin{tabular}{lll}
\hline
\textbf{Recurso} & \textbf{Estrutura} & \textbf{Mecanismo} \\
\hline
Ocupação das células do mapa & \texttt{Mapa.ocupada} & \texttt{Mapa.mutex} (global do mapa) \\
Estações de coleta (filas FIFO) & \texttt{Estacao.fila} & um \texttt{Estacao.mutex} por estação \\
Posições da esteira (\texttt{in}..\texttt{out}) & \texttt{Esteira.posicoes} & \texttt{Esteira.mutex} + \texttt{Esteira.cond} \\
Pontos de despacho \texttt{D} & células do mapa & via \texttt{Mapa.mutex} \\
Contadores/estatísticas & \texttt{Estatisticas} & \texttt{Estatisticas.mutex} \\
\textit{Flag} de encerramento & variável global & \texttt{mutex\_encerrar} \\
\hline
\end{tabular}
\end{table}

A ocupação do mapa é o recurso mais disputado: todo robô a lê e a escreve ao se
mover. As estações são filas FIFO acessadas pelo gerador (que enfileira) e pelos
coletores (que desenfileiram). A esteira é um \textit{buffer} limitado de uma
posição por célula. Os pontos \texttt{D} não têm trava própria: são células do
mapa e sua ocupação é coordenada pelo \texttt{Mapa.mutex}. Os contadores de
pacotes aguardando, na esteira e entregues, além do tempo de execução, são
atualizados por várias \textit{threads} e lidos pela interface.

\section{Seções Críticas}

As seções críticas identificadas e protegidas são:

\begin{itemize}
  \item \textbf{Ocupar uma célula (\textit{test-and-set}).} Verificar que a célula
        de destino está livre e marcá-la como ocupada precisa ser \emph{atômico};
        caso contrário, dois robôs poderiam ocupar a mesma célula. A verificação e
        a marcação ocorrem sob \texttt{Mapa.mutex}, e o passo do robô libera a
        célula de origem e ocupa a de destino dentro do mesmo \textit{lock},
        preservando a invariante física.
  \item \textbf{Fila da estação.} Enfileirar (gerador) e desenfileirar
        (coletores) ocorrem sob \texttt{Estacao.mutex}; vários coletores podem
        disputar a mesma fila simultaneamente.
  \item \textbf{Movimentação de pacotes na esteira.} Inserir em \texttt{in},
        avançar todas as posições (varrendo de \texttt{out} para \texttt{in} para
        não mover um pacote duas casas no mesmo passo) e retirar de \texttt{out}
        ocorrem sob \texttt{Esteira.mutex}; as consultas (saída ocupada, total de
        pacotes) também travam, garantindo leitura consistente.
  \item \textbf{Contadores globais.} Incrementos e decrementos dos pacotes
        aguardando, na esteira e entregues ocorrem sob \texttt{Estatisticas.mutex}.
  \item \textbf{Leitura da interface.} Cada parte mutável do quadro é copiada
        (\textit{snapshot}) sob o \textit{mutex} do recurso correspondente antes de
        ser desenhada, de modo que a tela nunca exibe um estado inconsistente.
\end{itemize}

\section{Mecanismos de Sincronização}

O mecanismo básico é a exclusão mútua com \texttt{pthread\_mutex\_t}, aplicada com
granularidade fina: um \textit{mutex} por recurso. A esteira mantém, além do
\textit{mutex}, uma variável de condição (\texttt{Esteira.cond}), sinalizada com
\texttt{pthread\_cond\_broadcast} a cada mudança de estado (inserção, avanço e
retirada), pronta para suportar espera bloqueante; no ritmo atual, as entidades
reavaliam o estado de forma temporizada (\texttt{usleep} por passo), sempre lendo
o estado protegido pelo respectivo \textit{mutex}. Cada
\texttt{pthread\_mutex\_init}/\texttt{pthread\_cond\_init} possui o
\texttt{\_destroy} correspondente antes do término (a \textit{flag} de encerramento
usa inicializador estático), e cada \texttt{pthread\_create} tem seu
\texttt{pthread\_join}, evitando vazamento de recursos e \textit{threads}
pendentes.

\section{Decisões de Projeto}

\begin{itemize}
  \item \textbf{Desenvolvimento incremental.} A lógica de cada entidade foi escrita
        como uma função de \emph{um passo} do seu autômato e validada de forma
        sequencial; a concorrência foi introduzida depois, reutilizando o mesmo
        passo como corpo do laço de cada \textit{thread}.
  \item \textbf{Granularidade fina de travas.} Um \textit{mutex} por recurso, em vez
        de um \textit{lock} global, para permitir progresso paralelo e reduzir
        contenção.
  \item \textbf{Navegação sem \textit{deadlock}.} Heurística gulosa (distância de
        Manhattan) com desvio lateral quando o caminho está bloqueado; robôs
        ociosos recuam a uma base para não obstruir \texttt{in}/\texttt{out},
        evitando o travamento de dois robôs frente a frente no corredor.
  \item \textbf{Memória previsível.} Pacotes e robôs vivem em vetores de tamanho
        fixo (\texttt{MAX\_*}), sem \texttt{malloc} por entidade.
  \item \textbf{\textit{Backpressure}.} Filas de coleta limitadas e gerador em
        \textit{round-robin} com \textit{liveness}, que procura outra estação com
        espaço em vez de desistir quando a primeira está cheia.
  \item \textbf{Verificação.} Testes unitários por módulo (\texttt{make test}) e
        detecção de condições de corrida com o \textit{ThreadSanitizer}
        (\texttt{make test-tsan}).
\end{itemize}

\section{Guia de Compilação e Execução}

Ambiente de referência: Linux (Ubuntu 24.04). Instalação das dependências
(compilador, \textit{make} e a biblioteca \texttt{ncurses}):

{\small
\begin{verbatim}
sudo apt install build-essential libncurses-dev
\end{verbatim}
}

Compilação (\texttt{gcc} com \texttt{-Wall -Wextra -pthread} e \texttt{-lncurses}),
execução e testes:

{\small
\begin{verbatim}
make                # gera o executável em build/idp
make run            # compila e executa (menu de cenários)
./build/idp 0       # executa diretamente o cenário 0, 1 ou 2
make test           # testes unitários dos módulos
make test-tsan      # teste de concorrência sob ThreadSanitizer
make clean          # remove artefatos de compilação
\end{verbatim}
}

Em terminal interativo, a interface \texttt{ncurses} exibe o mapa, os robôs, a
esteira e o painel (pacotes aguardando, na esteira, entregues e tempo de
execução); tecle \texttt{q} para sair. Sem terminal interativo, o programa imprime
um quadro final em texto com as mesmas informações, útil para validação
automatizada.

\begin{thebibliography}{9}
\bibitem{tanenbaum} TANENBAUM, A. S. \textit{Sistemas Operacionais Modernos}.
  3.~ed. São Paulo: Pearson Prentice Hall, 2010.
\bibitem{ncurses} FREE SOFTWARE FOUNDATION. \textit{ncurses --- A Free Software
  emulation of curses in System V Release 6.6}. Disponível em:
  \url{https://invisible-island.net/ncurses/announce.html}.
\end{thebibliography}

\end{document}
